\BourbakiExerciseGroup

\begin{enumerate}
\def\labelenumi{\arabic{enumi})}
\tightlist
\item
  \begin{enumerate}
  \def\labelenumii{\alph{enumii})}
  \tightlist
  \item
    试证：环 \(\mathbf{Z}_p\) （ \(p\) 进整数组成的环{[}§ 6 习题 23 a){]}）同构于离散环序列 \(\mathbf{Z}/(p^n)\) 的逆极限，其中 \(f_{nm}: \mathbf{Z}/(p^m) \to \mathbf{Z}/(p^n)\) （ \(n \leqslant m\) ）是把 \(\mathbf{Z}/(p^n)\) 看作 \(\mathbf{Z}/(p^m)\) 的商环时的典范同态。
  \end{enumerate}
\end{enumerate}

\begin{enumerate}
\def\labelenumi{\alph{enumi})}
\setcounter{enumi}{1}
\item
  对每个整数 \(n > 0\) ，令 \(G_n\) 为有理整数群 \(Z\) ，令 \(X_n\) 为其商群 \(Z/(p^n)\) （ \(p\) 素数），\(G_n\) 与 \(X_n\) 都赋有离散拓扑；把 \((G_n)\) 看作逆系统，\(G_m \to G_n\) 为恒等映射；并把 \((X_n)\) 看作逆系统，\(f_{nm}: X_m \to X_n\) 为典范同态（对 \(n \leqslant m\) ，见 a)）。则 \(X_n\) 的每点关于 \(G_n\) 的轨道是紧的，但点 \(x = (x_n)\) （ \(X = \varprojlim X_n\) 的）关于 \(G = \varprojlim G_n\) 的轨道不是紧的，并且不同构于诸 \(x_n\) 的轨道的逆极限；此外，典范映射 \(X/G \to \varprojlim X_n/G_n\) 不是单射的。c) 对每个 \(n > 0\) ，令 \(G_n\) 为子群 \(p^n Z\) （ \(Z\) 的，\(p\) 素数），令 \(X_n\) 为群 \(Z\) ，其中 \(G_n\) 与 \(X_n\) 赋有离散拓扑，且 \(G_n\) 以平移作用在 \(X_n\) 上。把 \((G_n)\) 与 \((X_n)\) 看作逆系统，\(X_m \to X_n\) 为恒等映射，\(G_m \to G_n\) 为典范单射（ \(n \leqslant m\) ）。则 \(X_n\) 的每点的稳定子群（关于 \(G_n\) 的）是紧的，\(X = \varprojlim X_n\) 的每点关于 \(G = \varprojlim G_n\) 的轨道是紧的，但典范映射 \(X/G = \varprojlim X_n/G_n\) 不是满射的。
\item
  利用 b) 与 c)，构造一个例子：带算子空间 \((\mathbf{X}_{\alpha})\) 关于群 \((\mathbf{G}_{\alpha})\) 的逆系统，使得若 \(X = \varprojlim X_{\alpha}\) 且 \(G = \varprojlim G_{\alpha}\) ，则典范映射 \(X/G \to \varprojlim X_{\alpha}/G_{\alpha}\) 既不是单射也不是满射。
\end{enumerate}

\begin{enumerate}
\def\labelenumi{\arabic{enumi})}
\setcounter{enumi}{1}
\tightlist
\item
  设 \((\mathbf{X}_{\alpha}, f_{\alpha \beta})\) 是非空集的一个逆系统，且 \(\varprojlim \mathbf{X}_{\alpha} = \emptyset\) 。令 \(\mathbf{G}_{\alpha}\) 表示以 \(\mathbf{Z}\) 为系数的 \(\mathbf{X}_{\alpha}\) 的元素的形式线性组合所成的自由 \(\mathbf{Z}\) -模，并令 \(g_{\alpha \beta}\) 表示同态 \(\mathbf{G}_{\beta} \to \mathbf{G}_{\alpha}\) （它与 \(f_{\alpha \beta}\) 在 \(\mathbf{X}_{\beta}\) 上一致）。试证：群 \(\mathbf{G} = \varprojlim \mathbf{G}_{\alpha}\) 仅由单位元素组成。{[}用反证法论证：若 \(z = (z_{\alpha})\) 是 \(\varprojlim \mathbf{G}_{\alpha}\) 的一个元素，则对每个 \(\alpha\) ，考察有限集 \(\mathbf{F}_{\alpha}\) （由 \(\mathbf{X}_{\alpha}\) 中在 \(z_{\alpha}\) 里的系数不为零的元素组成），并注意 \(f_{\alpha \beta}(\mathbf{F}_{\beta}) = \mathbf{F}_{\alpha}\) 。{]}由此构造群的一个逆系统 \((\mathbf{G}_{\alpha}, g_{\alpha \beta})\) 的例子，使得诸 \(g_{\alpha \beta}\) 是满射的，诸 \(\mathbf{G}_{\alpha}\) 是无限的，而 \(\varprojlim \mathbf{G}_{\alpha}\) 仅由单位元素组成。
\end{enumerate}

¶ 3) a) 试证每个全不连通紧群都是有限离散群族的一个逆极限（利用 § 4 习题 18）。b) 设 G 是一个全不连通紧群，L 是 G 的一个闭子群。试证存在与典范映射 G → G/L 相伴的连续截面 G/L → G。{[}利用 a) 与 § 3 命题 3；对每个 α 考察由截面 Gα/Lα → Gα 组成的有限集 Fα，并注意这些集关于典范满射 hαβ : Fβ → Fα 构成一个逆系统{]}。

\begin{enumerate}
\def\labelenumi{\arabic{enumi})}
\setcounter{enumi}{3}
\item
  设 \(G\) 是一个豪斯多夫拓扑群，\((H_{\alpha})\) 是 \(G\) 的紧正规子群的一个族，它关于关系 \(\supset\) 有向，且满足第 3 目的条件 (AP)。设 \((L_{\alpha})\) 是 \(G\) 的闭子群的一个族，使得 \(H_{\alpha} \subset L_{\alpha}\) 对每个 \(\alpha \in I\) 成立，且 \(L_{\alpha} = H_{\alpha}L_{\beta}\) 对 \(\alpha \leqslant \beta\) 成立，并令 \(L = \bigcap_{\alpha} L_{\alpha}\) 。试证 \(L_{\alpha} = H_{\alpha}L\) 对每个 \(\alpha\) 成立（利用命题 3）。
\item
  设 G 是一个豪斯多夫拓扑群，X 是 G 连续地作用其上的一个豪斯多夫拓扑空间；设 \((\mathrm{H}_{\alpha})\) 是 G 的紧正规子群的一个族，它关于关系 \(\supset\) 有向，且满足第 3 目的条件 (AP)。令 \(X_{\alpha}\) 表示 \(X/H_{\alpha}\) 。试证典范映射 \(X \to \varprojlim X_{\alpha}\) 是一个同胚。
\end{enumerate}

¶ * 6) 设 \((\mathbf{G}_n, f_{nm})_{n \in \mathbb{N}}\) 是紧群的一个逆系统，使得 \(\mathbf{G}_n = \mathbf{T} = \mathbf{R} / \mathbf{Z}\) 对每个 \(n\) 成立，且 \(f_{nm}\) 是连续同态 \(x \to p^{m-n}x\) ，它把 \(\mathbf{T}\) 映到其自身上（对 \(n \leqslant m\) ；\(p\) 为给定的素数）。拓扑群 \(\mathbf{T}_p = \varprojlim \mathbf{G}_n\) 称为 \(p\) 进螺线管；它是一个紧连通交换拓扑群。

\begin{enumerate}
\def\labelenumi{\alph{enumi})}
\item
  对每个 n，连续同态 \(f_{n}: T_{p} \to G_{n}\) 是满射的，其核同构于 p 进整数群 \(Z_{p}\) {[}参看习题 1 a){]}。
\item
  设 \(\varphi\) 是典范同态 \(\mathbf{R} \to \mathbf{R} / \mathbf{Z} = \mathbf{T}\) 。对每个 \(x \in \mathbf{R}\) ，令 \(\theta(x) = (\varphi(x / p^n))_{n \in \mathbb{N}}\) ；试证 \(\theta\) 是 \(\mathbf{R}\) 到 \(\mathbf{T}_p\) 中的单射连续同态，且 \(\theta(\mathbf{R})\) 是 \(\mathbf{T}_p\) 的一个稠密子群。
\item
  设 \(\mathbf{I}\) 是 \(\mathbf{R}\) 中的一个开区间，以 \(0\) 为中心且长度 \(< 1\) 。试证子空间 \(\overline{f}_{0}^{1}(\varphi(\mathbf{I}))\) （ \(\mathbf{T}_p\) 的）同胚于积 \(\mathbf{I} \times \mathbf{Z}_p\) 。特别地，群 \(\mathbf{T}_p\) 不是局部连通的。
\item
  试证每个闭子群 \(\mathbf{H}\) （ \(\mathbf{T}_p\) 的，异于 \(\mathbf{T}_p\) 与 \(\{0\}\) ）都全不连通，且同构于形如 \(\mathbf{Z}/n\mathbf{Z}\) 或 \((\mathbf{Z}/n\mathbf{Z}) \times \mathbf{Z}_p\) 的群，其中 \(n\) 是与 \(p\) 互素的整数（利用第 3 目命题 3）。
\item
  试证 \(\mathbf{T}_p\) 是一个不可分解的紧连通空间，即 \(\mathbf{T}_p\) 不能被两个紧连通集 \(\mathbf{P}\) 、\(\mathbf{Q}\) （都不等于 \(\mathbf{T}_p\) ）所覆盖。（注意存在整数 \(n\) ，使得 \(f_n(\mathbf{P}) \neq \mathbf{G}_n\) 且 \(f_n(\mathbf{Q}) \neq \mathbf{G}_n\) ，并考察集 \(f_{n+1}(\mathbf{P})\) 与 \(f_{n+1}(\mathbf{Q})\) 以得出矛盾）。*
\end{enumerate}

